\documentclass[10pt,a4paper]{amsart}
\usepackage{amsmath,amssymb,amsthm,graphicx}
\usepackage{verbatim}
\newenvironment{reference}{\comment}{\endcomment}
%\newenvironment{reference}{}{}
\newenvironment{slogan}{\comment}{\endcomment}
\newenvironment{history}{\comment}{\endcomment}
% For commutative diagrams we use Xy-pic
\usepackage[all]{xy}

% We use 2cell for 2-commutative diagrams.
\xyoption{2cell}
% Theorem environments.
%
\theoremstyle{plain}
\newtheorem{theorem}[subsection]{Theorem}
\newtheorem{proposition}[subsection]{Proposition}
\newtheorem{lemma}[subsection]{Lemma}

\theoremstyle{definition}
\newtheorem{definition}[subsection]{Definition}
\newtheorem{example}[subsection]{Example}
\newtheorem{exercise}[subsection]{Exercise}
\newtheorem{situation}[subsection]{Situation}

\theoremstyle{remark}
\newtheorem{remark}[subsection]{Remark}
\newtheorem{remarks}[subsection]{Remarks}

\numberwithin{equation}{subsection}

% Macros
%
\def\lim{\mathop{\mathrm{lim}}\nolimits}
\def\colim{\mathop{\mathrm{colim}}\nolimits}
\def\Spec{\mathop{\mathrm{Spec}}}
\def\Hom{\mathop{\mathrm{Hom}}\nolimits}
\def\Ext{\mathop{\mathrm{Ext}}\nolimits}
\def\SheafHom{\mathop{\mathcal{H}\!\mathit{om}}\nolimits}
\def\SheafExt{\mathop{\mathcal{E}\!\mathit{xt}}\nolimits}
\def\Sch{\mathit{Sch}}
\def\Mor{\mathop{\mathrm{Mor}}\nolimits}
\def\Ob{\mathop{\mathrm{Ob}}\nolimits}
\def\Sh{\mathop{\mathit{Sh}}\nolimits}
\def\NL{\mathop{N\!L}\nolimits}
\def\CH{\mathop{\mathrm{CH}}\nolimits}
\def\proetale{{pro\text{-}\acute{e}tale}}
\def\etale{{\acute{e}tale}}
\def\QCoh{\mathit{QCoh}}
\def\Ker{\mathop{\mathrm{Ker}}}
\def\Im{\mathop{\mathrm{Im}}}
\def\Coker{\mathop{\mathrm{Coker}}}
\def\Coim{\mathop{\mathrm{Coim}}}

% Boxtimes
%
\DeclareMathSymbol{\boxtimes}{\mathbin}{AMSa}{"02}

%
% Macros for moduli stacks/spaces
%
\def\QCohstack{\mathcal{QC}\!\mathit{oh}}
\def\Cohstack{\mathcal{C}\!\mathit{oh}}
\def\Spacesstack{\mathcal{S}\!\mathit{paces}}
\def\Quotfunctor{\mathrm{Quot}}
\def\Hilbfunctor{\mathrm{Hilb}}
\def\Curvesstack{\mathcal{C}\!\mathit{urves}}
\def\Polarizedstack{\mathcal{P}\!\mathit{olarized}}
\def\Complexesstack{\mathcal{C}\!\mathit{omplexes}}
% \Pic is the operator that assigns to X its picard group, usage \Pic(X)
% \Picardstack_{X/B} denotes the Picard stack of X over B
% \Picardfunctor_{X/B} denotes the Picard functor of X over B
\def\Pic{\mathop{\mathrm{Pic}}\nolimits}
\def\Picardstack{\mathcal{P}\!\mathit{ic}}
\def\Picardfunctor{\mathrm{Pic}}
\def\Deformationcategory{\mathcal{D}\!\mathit{ef}}
\usepackage{cleveref}
\crefname{originallemma}{original lemma}{original lemmas}
\Crefname{originallemma}{Original lemma}{Original lemmas}
\makeatletter
\newlabel{lemma-fibre-product-after-map}{{lemma-fibre-product-after-map}{source}{}{sourcepointer-1}{}}
\newlabel{lemma-fibre-product-after-map@cref}{{[originallemma][0][]lemma-fibre-product-after-map}{[1][0][]source}{}{}{}}
\makeatother
\crefname{originallemma}{original lemma}{original lemmas}
\Crefname{originallemma}{Original lemma}{Original lemmas}
\makeatletter
\newlabel{lemma-properties-diagonal}{{lemma-properties-diagonal}{source}{}{sourcepointer-2}{}}
\newlabel{lemma-properties-diagonal@cref}{{[originallemma][0][]lemma-properties-diagonal}{[1][0][]source}{}{}{}}
\makeatother
\crefname{originalreference}{original reference}{original references}
\Crefname{originalreference}{Original reference}{Original references}
\makeatletter
\newlabel{spaces-morphisms-lemma-base-change-quasi-finite}{{spaces-morphisms-lemma-base-change-quasi-finite}{source}{}{sourcepointer-3}{}}
\newlabel{spaces-morphisms-lemma-base-change-quasi-finite@cref}{{[originalreference][0][]spaces-morphisms-lemma-base-change-quasi-finite}{[1][0][]source}{}{}{}}
\makeatother
\crefname{originalreference}{original reference}{original references}
\Crefname{originalreference}{Original reference}{Original references}
\makeatletter
\newlabel{spaces-morphisms-lemma-base-change-unramified}{{spaces-morphisms-lemma-base-change-unramified}{source}{}{sourcepointer-4}{}}
\newlabel{spaces-morphisms-lemma-base-change-unramified@cref}{{[originalreference][0][]spaces-morphisms-lemma-base-change-unramified}{[1][0][]source}{}{}{}}
\makeatother
\crefname{originalreference}{original reference}{original references}
\Crefname{originalreference}{Original reference}{Original references}
\makeatletter
\newlabel{spaces-morphisms-lemma-compose-after-separated}{{spaces-morphisms-lemma-compose-after-separated}{source}{}{sourcepointer-5}{}}
\newlabel{spaces-morphisms-lemma-compose-after-separated@cref}{{[originalreference][0][]spaces-morphisms-lemma-compose-after-separated}{[1][0][]source}{}{}{}}
\makeatother
\crefname{originalreference}{original reference}{original references}
\Crefname{originalreference}{Original reference}{Original references}
\makeatletter
\newlabel{spaces-morphisms-lemma-permanence-quasi-finite}{{spaces-morphisms-lemma-permanence-quasi-finite}{source}{}{sourcepointer-6}{}}
\newlabel{spaces-morphisms-lemma-permanence-quasi-finite@cref}{{[originalreference][0][]spaces-morphisms-lemma-permanence-quasi-finite}{[1][0][]source}{}{}{}}
\makeatother
\crefname{originalreference}{original reference}{original references}
\Crefname{originalreference}{Original reference}{Original references}
\makeatletter
\newlabel{spaces-morphisms-lemma-permanence-unramified}{{spaces-morphisms-lemma-permanence-unramified}{source}{}{sourcepointer-7}{}}
\newlabel{spaces-morphisms-lemma-permanence-unramified@cref}{{[originalreference][0][]spaces-morphisms-lemma-permanence-unramified}{[1][0][]source}{}{}{}}
\makeatother
\crefname{originalreference}{original reference}{original references}
\Crefname{originalreference}{Original reference}{Original references}
\makeatletter
\newlabel{spaces-morphisms-lemma-quasi-compact-permanence}{{spaces-morphisms-lemma-quasi-compact-permanence}{source}{}{sourcepointer-8}{}}
\newlabel{spaces-morphisms-lemma-quasi-compact-permanence@cref}{{[originalreference][0][]spaces-morphisms-lemma-quasi-compact-permanence}{[1][0][]source}{}{}{}}
\makeatother
\crefname{originalreference}{original reference}{original references}
\Crefname{originalreference}{Original reference}{Original references}
\makeatletter
\newlabel{spaces-morphisms-lemma-universally-closed-permanence}{{spaces-morphisms-lemma-universally-closed-permanence}{source}{}{sourcepointer-9}{}}
\newlabel{spaces-morphisms-lemma-universally-closed-permanence@cref}{{[originalreference][0][]spaces-morphisms-lemma-universally-closed-permanence}{[1][0][]source}{}{}{}}
\makeatother
\title{The first factor of a DM, quasi-DM, separated or quasi-separated composite has the same property}
\author{The Stacks Project Authors}
\date{stacks-project:stacks-morphisms.tex:local-snapshot-20261012; GNU FDL 1.2 or later}
\begin{document}
\maketitle
\noindent\textit{Editorial note.} All delivered mathematical text is the author's own, copied from the pinned Stacks Project chapter source sources/stacks-algebraic-stacks-morphisms/pinned-original-source.tex; no statement or proof step was rewritten, reordered or removed.\par\medskip
\noindent\textit{Editorial note.} Preamble provenance: the Stacks Project keeps one single shared preamble for every chapter; this card takes it directly from the staged sources/stacks-algebraic-stacks-morphisms/preamble.tex as three verbatim spans (lines 8-12, 14-18 and 154-217), so no derived file was written and nothing was added or redefined. The wrapper class and the base packages are supplied by the assembler.\par\medskip
\noindent\textit{Editorial note.} Bibliography provenance: neither the counted statement nor its proof contains a citation (this lane rescans the whole span before writing the card), so the card carries no bibliography and no entry was reconstructed.\par\medskip
\noindent\textit{Editorial note.} Selection provenance: this card was opened under the file-based control contract (CONTROL.md), with the mechanical pre-card gate run on its statement and proof interval before the card was written.\par\medskip
\noindent\textit{Allowed standard prerequisites (editorial).}\begin{itemize}
\item The characterisation of DM, quasi-DM, separated and quasi-separated morphisms of algebraic stacks through properties of the diagonal, and the representability of the morphisms in the factorisation of a diagonal. Invoked at the start of the counted proof (source lines 949-952 and again in the other parts) as pointers to two lemmas of this chapter. Turns the hypotheses into properties of the induced morphisms of algebraic spaces.
\item The permanence and base change statements for unramified, locally quasi-finite, quasi-separated, proper and quasi-compact morphisms of algebraic spaces. Invoked repeatedly inside the counted proof (source lines 964-989, 1021-1026, 1054-1056) as cross-chapter pointers. Transfers the property from the composite of the comparison morphisms to the intermediate one and then to the diagonal of the morphism in question.
\end{itemize}\medskip
\ifdefined\hypertarget\hypertarget{sourcepointer-1}{}\fi
\noindent\textit{Source reference.} Original lemma lemma-fibre-product-after-map: sources/stacks-algebraic-stacks-morphisms/pinned-original-source.tex, source lines 824--824. Navigation pointer only: the counted text refers to the Stacks item lemma-fibre-product-after-map, which belongs to another chapter of the project and is not part of this card. The pointer carries the target's own label as its number because this chapter's staged snapshot does not contain the other chapter's tag table; its content is not reproduced here.\par\medskip
\ifdefined\hypertarget\hypertarget{sourcepointer-2}{}\fi
\noindent\textit{Source reference.} Original lemma lemma-properties-diagonal: sources/stacks-algebraic-stacks-morphisms/pinned-original-source.tex, source lines 420--420. Navigation pointer only: the counted text refers to the Stacks item lemma-properties-diagonal, which belongs to another chapter of the project and is not part of this card. The pointer carries the target's own label as its number because this chapter's staged snapshot does not contain the other chapter's tag table; its content is not reproduced here.\par\medskip
\ifdefined\hypertarget\hypertarget{sourcepointer-3}{}\fi
\noindent\textit{Source reference.} Original reference spaces-morphisms-lemma-base-change-quasi-finite: sources/stacks-algebraic-stacks-morphisms/pinned-original-source.tex, source lines 989--989. Navigation pointer only: the counted text refers to the Stacks item spaces-morphisms-lemma-base-change-quasi-finite, which belongs to another chapter of the project and is not part of this card. The pointer carries the target's own label as its number because this chapter's staged snapshot does not contain the other chapter's tag table; its content is not reproduced here.\par\medskip
\ifdefined\hypertarget\hypertarget{sourcepointer-4}{}\fi
\noindent\textit{Source reference.} Original reference spaces-morphisms-lemma-base-change-unramified: sources/stacks-algebraic-stacks-morphisms/pinned-original-source.tex, source lines 986--986. Navigation pointer only: the counted text refers to the Stacks item spaces-morphisms-lemma-base-change-unramified, which belongs to another chapter of the project and is not part of this card. The pointer carries the target's own label as its number because this chapter's staged snapshot does not contain the other chapter's tag table; its content is not reproduced here.\par\medskip
\ifdefined\hypertarget\hypertarget{sourcepointer-5}{}\fi
\noindent\textit{Source reference.} Original reference spaces-morphisms-lemma-compose-after-separated: sources/stacks-algebraic-stacks-morphisms/pinned-original-source.tex, source lines 926--926. Navigation pointer only: the counted text refers to the Stacks item spaces-morphisms-lemma-compose-after-separated, which belongs to another chapter of the project and is not part of this card. The pointer carries the target's own label as its number because this chapter's staged snapshot does not contain the other chapter's tag table; its content is not reproduced here.\par\medskip
\ifdefined\hypertarget\hypertarget{sourcepointer-6}{}\fi
\noindent\textit{Source reference.} Original reference spaces-morphisms-lemma-permanence-quasi-finite: sources/stacks-algebraic-stacks-morphisms/pinned-original-source.tex, source lines 971--971. Navigation pointer only: the counted text refers to the Stacks item spaces-morphisms-lemma-permanence-quasi-finite, which belongs to another chapter of the project and is not part of this card. The pointer carries the target's own label as its number because this chapter's staged snapshot does not contain the other chapter's tag table; its content is not reproduced here.\par\medskip
\ifdefined\hypertarget\hypertarget{sourcepointer-7}{}\fi
\noindent\textit{Source reference.} Original reference spaces-morphisms-lemma-permanence-unramified: sources/stacks-algebraic-stacks-morphisms/pinned-original-source.tex, source lines 968--968. Navigation pointer only: the counted text refers to the Stacks item spaces-morphisms-lemma-permanence-unramified, which belongs to another chapter of the project and is not part of this card. The pointer carries the target's own label as its number because this chapter's staged snapshot does not contain the other chapter's tag table; its content is not reproduced here.\par\medskip
\ifdefined\hypertarget\hypertarget{sourcepointer-8}{}\fi
\noindent\textit{Source reference.} Original reference spaces-morphisms-lemma-quasi-compact-permanence: sources/stacks-algebraic-stacks-morphisms/pinned-original-source.tex, source lines 1026--1026. Navigation pointer only: the counted text refers to the Stacks item spaces-morphisms-lemma-quasi-compact-permanence, which belongs to another chapter of the project and is not part of this card. The pointer carries the target's own label as its number because this chapter's staged snapshot does not contain the other chapter's tag table; its content is not reproduced here.\par\medskip
\ifdefined\hypertarget\hypertarget{sourcepointer-9}{}\fi
\noindent\textit{Source reference.} Original reference spaces-morphisms-lemma-universally-closed-permanence: sources/stacks-algebraic-stacks-morphisms/pinned-original-source.tex, source lines 1056--1056. Navigation pointer only: the counted text refers to the Stacks item spaces-morphisms-lemma-universally-closed-permanence, which belongs to another chapter of the project and is not part of this card. The pointer carries the target's own label as its number because this chapter's staged snapshot does not contain the other chapter's tag table; its content is not reproduced here.\par\medskip
\begin{lemma}
\label{lemma-compose-after-separated}
Let $f : \mathcal{X} \to \mathcal{Y}$ and
$g : \mathcal{Y} \to \mathcal{Z}$ be morphisms of algebraic stacks.
\begin{enumerate}
\item If $g \circ f$ is DM then so is $f$.
\item If $g \circ f$ is quasi-DM then so is $f$.
\item If $g \circ f$ is separated and $\Delta_g$ is separated, then
$f$ is separated.
\item If $g \circ f$ is quasi-separated and
$\Delta_g$ is quasi-separated, then $f$ is quasi-separated.
\end{enumerate}
\end{lemma}

\begin{proof}
Consider the factorization
$$
\mathcal{X} \to
\mathcal{X} \times_\mathcal{Y} \mathcal{X} \to
\mathcal{X} \times_\mathcal{Z} \mathcal{X}
$$
of the diagonal morphism of $g \circ f$. Both morphisms are representable by
algebraic spaces, see
Lemmas \ref{lemma-properties-diagonal} and
\ref{lemma-fibre-product-after-map}.
Hence for any scheme $T$ and morphism
$T \to \mathcal{X} \times_\mathcal{Y} \mathcal{X}$
we get morphisms of algebraic spaces
$$
A = \mathcal{X} \times_{(\mathcal{X} \times_\mathcal{Z} \mathcal{X})} T
\longrightarrow
B = (\mathcal{X} \times_\mathcal{Y} \mathcal{X})
\times_{(\mathcal{X} \times_\mathcal{Z} \mathcal{X})} T
\longrightarrow
T
$$
If $g \circ f$ is DM (resp.\ quasi-DM), then the composition $A \to T$
is unramified (resp.\ locally quasi-finite). Hence $A \to B$ is
unramified (resp.\ locally quasi-finite) by
Morphisms of Spaces, Lemma
\ref{spaces-morphisms-lemma-permanence-unramified}
(resp.
Morphisms of Spaces,
Lemma \ref{spaces-morphisms-lemma-permanence-quasi-finite}).
Now consider the diagram
$$
\xymatrix{
A \times_B T \ar[r] \ar[d] & T \ar[d] \\
A \ar[r] \ar[d] & B \ar[r] \ar[d] & T \ar[d] \\
\mathcal{X} \ar[r] &
\mathcal{X} \times_\mathcal{Y} \mathcal{X} \ar[r] &
\mathcal{X} \times_\mathcal{Z} \mathcal{X}
}
$$
with all squares $2$-cartesian where the arrow $T \to B$ comes from the arrow
$T \to \mathcal{X} \times_\mathcal{Y} \mathcal{X}$ we started with.
The arrow $A \times_B T \to T$ is unramified (resp.\ locally quasi-finite)
by Morphisms of Spaces, Lemma
\ref{spaces-morphisms-lemma-base-change-unramified}
(resp.
Morphisms of Spaces,
Lemma \ref{spaces-morphisms-lemma-base-change-quasi-finite}).
Noting that $A \times_B T$ is equal to
$\mathcal{X} \times_{\mathcal{X} \times_\mathcal{Y} \mathcal{X}} T$
we conclude that $f$ is DM (resp.\ quasi-DM).
This proves (1) and (2).

\medskip\noindent
Proof of (4). Assume $g \circ f$ is quasi-separated and $\Delta_g$ is
quasi-separated. Consider the factorization
$$
\mathcal{X} \to
\mathcal{X} \times_\mathcal{Y} \mathcal{X} \to
\mathcal{X} \times_\mathcal{Z} \mathcal{X}
$$
of the diagonal morphism of $g \circ f$. Both morphisms are
representable by algebraic spaces and the second one is quasi-separated, see
Lemmas \ref{lemma-properties-diagonal} and
\ref{lemma-fibre-product-after-map}.
Hence for any scheme $T$ and morphism
$T \to \mathcal{X} \times_\mathcal{Y} \mathcal{X}$
we get morphisms of algebraic spaces
$$
A = \mathcal{X} \times_{(\mathcal{X} \times_\mathcal{Z} \mathcal{X})} T
\longrightarrow
B = (\mathcal{X} \times_\mathcal{Y} \mathcal{X})
\times_{(\mathcal{X} \times_\mathcal{Z} \mathcal{X})} T
\longrightarrow
T
$$
such that $B \to T$ is quasi-separated.
The composition $A \to T$ is quasi-compact and quasi-separated
as we have assumed that $g \circ f$ is quasi-separated.
Hence $A \to B$ is quasi-separated by
Morphisms of Spaces,
Lemma \ref{spaces-morphisms-lemma-compose-after-separated}.
And $A \to B$ is quasi-compact by
Morphisms of Spaces,
Lemma \ref{spaces-morphisms-lemma-quasi-compact-permanence}.
Thus $f$ is quasi-separated.

\medskip\noindent
Proof of (3). Assume $g \circ f$ is separated and $\Delta_g$ is
separated. Consider the factorization
$$
\mathcal{X} \to
\mathcal{X} \times_\mathcal{Y} \mathcal{X} \to
\mathcal{X} \times_\mathcal{Z} \mathcal{X}
$$
of the diagonal morphism of $g \circ f$. Both morphisms are
representable by algebraic spaces and the second one is separated, see
Lemmas \ref{lemma-properties-diagonal} and
\ref{lemma-fibre-product-after-map}.
Hence for any scheme $T$ and morphism
$T \to \mathcal{X} \times_\mathcal{Y} \mathcal{X}$
we get morphisms of algebraic spaces
$$
A = \mathcal{X} \times_{(\mathcal{X} \times_\mathcal{Z} \mathcal{X})} T
\longrightarrow
B = (\mathcal{X} \times_\mathcal{Y} \mathcal{X})
\times_{(\mathcal{X} \times_\mathcal{Z} \mathcal{X})} T
\longrightarrow
T
$$
such that $B \to T$ is separated.
The composition $A \to T$ is proper as we have assumed that
$g \circ f$ is quasi-separated. Hence $A \to B$ is proper by
Morphisms of Spaces,
Lemma \ref{spaces-morphisms-lemma-universally-closed-permanence}
which means that $f$ is separated.
\end{proof}

\end{document}
